% results.tex -- BNR-2026-012

\section{Results}
\label{sec:results}

\subsection{Attestation Pipeline Latency}

Table~\ref{tab:attestation} presents the attestation pipeline latency
across five batch sizes with 50 replications per configuration.

\begin{table}[t]
\centering
\caption{Attestation Pipeline Latency (ms), $(5,7)$ Configuration}
\label{tab:attestation}
\begin{tabular}{@{}lrrrrr@{}}
\toprule
\textbf{Batch} & \textbf{Mean} & \textbf{P50} & \textbf{P95} & \textbf{CI\textsubscript{95}} & \textbf{Margin} \\
\midrule
10\,KB  & 0.292  & 0.100  & 0.600  & [0.21, 0.37]  & 1{,}667$\times$ \\
100\,KB & 1.121  & 0.942  & 2.051  & [0.96, 1.28]  & 488$\times$     \\
500\,KB & 4.509  & 4.370  & 5.393  & [4.34, 4.68]  & 185$\times$     \\
1\,MB   & 8.940  & 8.692  & 10.375 & [8.68, 9.20]  & 96$\times$      \\
5\,MB   & 46.871 & 46.020 & 52.675 & [45.96, 47.79] & 19$\times$     \\
\bottomrule
\multicolumn{6}{l}{\footnotesize Target: $<$1{,}000\,ms. All configurations PASS.} \\
\end{tabular}
\end{table}

Attestation latency scales linearly at approximately 9\,$\mu$s/KB
(0.009\,ms/KB). All configurations pass the $<$1 second target with
margins between $19\times$ (5\,MB) and $1{,}667\times$ (10\,KB). The
Go implementation with SIMD-optimized RS encoding achieves a
$36\times$ speedup over the TypeScript BigInt-based RU-V6 baseline:
163.5\,ms $\rightarrow$ 4.5\,ms at 500\,KB and 320\,ms $\rightarrow$
8.9\,ms at 1\,MB.

\subsection{Storage Overhead}

Table~\ref{tab:storage} presents per-node and total storage across batch
sizes.

\begin{table}[t]
\centering
\caption{Storage Overhead, $(5,7)$ RS Configuration}
\label{tab:storage}
\begin{tabular}{@{}lrrrr@{}}
\toprule
\textbf{Batch} & \textbf{Per-Node} & \textbf{Total (7)} & \textbf{Ratio} & \textbf{RU-V6} \\
\midrule
10\,KB  & 2\,KB     & 14\,KB    & 1.47$\times$ & 3.87$\times$ \\
100\,KB & 20\,KB    & 140\,KB   & 1.41$\times$ & 3.87$\times$ \\
500\,KB & 100\,KB   & 700\,KB   & 1.40$\times$ & 3.87$\times$ \\
1\,MB   & 204\,KB   & 1{,}434\,KB & 1.40$\times$ & 3.87$\times$ \\
5\,MB   & 1{,}024\,KB & 7{,}168\,KB & 1.40$\times$ & 3.87$\times$ \\
\bottomrule
\multicolumn{5}{l}{\footnotesize Theoretical minimum for $(5,7)$ RS: $7/5 = 1.40\times$.} \\
\end{tabular}
\end{table}

Storage overhead converges to $1.40\times$ for batch sizes $\geq$
500\,KB, matching the theoretical minimum of $n/k = 7/5$. This
represents a $2.77\times$ improvement over RU-V6 Shamir ($3.87\times$).
Per-node storage is $1/5$ of the original data, compared to $1\times$
(full copy) for Shamir. For a 1\,MB batch, each node stores 204\,KB
versus 1{,}290\,KB with Shamir---a $6.3\times$ per-node improvement.

\subsection{Recovery Time}

Table~\ref{tab:recovery} presents recovery benchmarks with full and
degraded committees.

\begin{table}[t]
\centering
\caption{Recovery Time (ms), 30 Replications}
\label{tab:recovery}
\begin{tabular}{@{}llrrrr@{}}
\toprule
\textbf{Batch} & \textbf{Nodes} & \textbf{Mean} & \textbf{P95} & \textbf{OK} & \textbf{Match} \\
\midrule
10\,KB  & 7 (all)   & 0.067 & 0.533 & 100\% & 100\% \\
100\,KB & 7 (all)   & 0.141 & 1.006 & 100\% & 100\% \\
500\,KB & 7 (all)   & 0.677 & 1.747 & 100\% & 100\% \\
1\,MB   & 7 (all)   & 0.949 & 1.959 & 100\% & 100\% \\
\midrule
100\,KB & 5 (2 down) & 0.103 & ---   & 100\% & 30/30 \\
500\,KB & 5 (2 down) & 0.418 & ---   & 100\% & 30/30 \\
1\,MB   & 5 (2 down) & 0.665 & ---   & 100\% & 30/30 \\
\bottomrule
\end{tabular}
\end{table}

Recovery at 1\,MB completes in under 1\,ms, compared to 2{,}482\,ms for
RU-V6 Lagrange interpolation---a $2{,}613\times$ speedup. Recovery with
2 nodes down is actually faster than with all nodes available (fewer
shards to process). All recovered data matches the original byte-for-byte.

\subsection{Failure Tolerance}

Table~\ref{tab:failure} validates correct behavior under progressive node
failures.

\begin{table}[t]
\centering
\caption{Failure Tolerance Validation (500\,KB batch, 30 reps)}
\label{tab:failure}
\begin{tabular}{@{}lccccl@{}}
\toprule
\textbf{Down} & \textbf{Online} & \textbf{Disperse} & \textbf{Recover} & \textbf{Match} & \textbf{Expected} \\
\midrule
0 & 7 & 100\% & 30/30 & 30/30 & PASS \\
1 & 6 & 100\% & 30/30 & 30/30 & PASS \\
2 & 5 & 100\% & 30/30 & 30/30 & PASS \\
3 & 4 & 0\%   & 0/30  & 0/30  & PASS \\
\bottomrule
\multicolumn{6}{l}{\footnotesize 3 nodes down correctly triggers fallback (INV-DA-P5).} \\
\end{tabular}
\end{table}

The system correctly tolerates 0, 1, and 2 node failures with 100\%
dispersal and recovery success. At 3 nodes down (4 online, below the
threshold of 5), dispersal correctly fails, triggering the on-chain DA
fallback path as specified by INV-DA-P5.

\subsection{Privacy Validation}

Table~\ref{tab:privacy} presents the privacy test results.

\begin{table}[t]
\centering
\caption{Privacy Validation Results (30 Replications)}
\label{tab:privacy}
\begin{tabular}{@{}lrl@{}}
\toprule
\textbf{Test} & \textbf{Result} & \textbf{Details} \\
\midrule
Chunk independence     & PASS   & Different batches $\rightarrow$ different chunks \\
$k{-}1$ key secrecy   & 30/30  & 4 shares never recover key \\
$k$ shares recovery   & 30/30  & 5 shares always recover key \\
Round-trip integrity   & 40/40  & Byte-perfect at 1--500\,KB \\
Certificate verification & 30/30 & ECDSA signatures verify \\
\bottomrule
\end{tabular}
\end{table}

All privacy tests pass with zero failures. The $k - 1 = 4$ share test
confirms that fewer than 5 Shamir shares never recover the AES key
(information-theoretic guarantee per INV-DA-P3). The round-trip
integrity test validates the complete pipeline: encrypt $\rightarrow$
RS encode $\rightarrow$ distribute $\rightarrow$ collect $\rightarrow$
RS decode $\rightarrow$ Shamir recover $\rightarrow$ decrypt, with
byte-perfect reconstruction across batch sizes from 1\,KB to 500\,KB.

\subsection{Configuration Comparison: (5,7) vs (2,3)}

Table~\ref{tab:comparison} compares the proposed $(5,7)$ configuration
against a $(2,3)$ baseline using the same hybrid protocol.

\begin{table}[t]
\centering
\caption{Configuration Comparison: $(5,7)$ vs $(2,3)$, 50 Replications}
\label{tab:comparison}
\begin{tabular}{@{}lrrrrr@{}}
\toprule
\textbf{Batch} & \textbf{(5,7)} & \textbf{(2,3)} & \textbf{Ratio} & \textbf{Stor.\,(5,7)} & \textbf{Stor.\,(2,3)} \\
\midrule
100\,KB & 1.105 & 1.002 & 1.10$\times$ & 1.41$\times$ & 1.50$\times$ \\
500\,KB & 4.525 & 4.762 & 0.95$\times$ & 1.40$\times$ & 1.50$\times$ \\
1\,MB   & 9.490 & 9.000 & 1.05$\times$ & 1.40$\times$ & 1.50$\times$ \\
\bottomrule
\multicolumn{6}{l}{\footnotesize Latency in ms. Stor. = total storage overhead ratio.} \\
\end{tabular}
\end{table}

The latency overhead of the 7-node configuration versus 3-node is
negligible ($1.0$--$1.1\times$). At 500\,KB, the $(5,7)$ configuration
is actually slightly faster ($0.95\times$), because RS encoding is
$O(n \log n)$ regardless of $k$ and the AES encryption overhead is
identical. The $(5,7)$ configuration achieves lower total storage
overhead ($1.40\times$ vs $1.50\times$) due to the more favorable
$n/k$ ratio.

\subsection{Benchmark Reconciliation}

Table~\ref{tab:reconciliation} validates experimental results against
published benchmarks.

\begin{table}[t]
\centering
\caption{Benchmark Reconciliation with Published Results}
\label{tab:reconciliation}
\begin{tabular}{@{}p{2.5cm}rrp{1.8cm}@{}}
\toprule
\textbf{Metric} & \textbf{Literature} & \textbf{Measured} & \textbf{Assessment} \\
\midrule
EigenDA V2 E2E & 5\,s (200+\,nodes) & 4.5\,ms (7\,nodes) & Expected \\
RU-V6 att. 500\,KB & 163.5\,ms & 4.5\,ms & 36$\times$ faster \\
RU-V6 att. 1\,MB & 320\,ms & 8.9\,ms & 36$\times$ faster \\
RU-V6 rec. 1\,MB & 2{,}482\,ms & 0.95\,ms & 2{,}613$\times$ \\
RU-V6 storage & 3.87$\times$ & 1.40$\times$ & 2.77$\times$ better \\
AES-GCM 1\,MB & $\sim$0.1\,ms & Included & Consistent \\
Shamir 32\,B & $\sim$17\,$\mu$s & Included & Negligible \\
\bottomrule
\end{tabular}
\end{table}

No result diverges by more than $10\times$ from literature estimates for
comparable configurations. The $36\times$ attestation speedup over RU-V6
is attributable to three factors: compiled Go versus interpreted
JavaScript, SIMD-optimized RS versus BigInt Shamir field operations,
and hardware AES-NI acceleration. The $1{,}111\times$ difference from
EigenDA V2 is expected given the difference in committee size (7 versus
200+ nodes) and the absence of network latency in our local benchmarks.
